% Eldrun — Detailed README (LaTeX edition)
% Compile with: pdflatex README.tex  (run twice for the table of contents)
\documentclass[11pt,a4paper]{article}

\usepackage[utf8]{inputenc}
\usepackage[T1]{fontenc}
\usepackage{lmodern}
\usepackage[margin=2.4cm]{geometry}
\usepackage{microtype}
\usepackage{xcolor}
\usepackage{booktabs}
\usepackage{longtable}
\usepackage{array}
\usepackage{enumitem}
\usepackage{titlesec}
\usepackage{fancyhdr}
\usepackage{listings}
\usepackage[hidelinks,colorlinks=true,linkcolor=eldrunblue,urlcolor=eldrunblue]{hyperref}

% ---- Colours ---------------------------------------------------------------
\definecolor{eldrunblue}{HTML}{2D6CDF}
\definecolor{eldrungray}{HTML}{4A4A4A}
\definecolor{codebg}{HTML}{F4F5F7}
\definecolor{donegreen}{HTML}{1B7F3B}
\definecolor{openorange}{HTML}{B5651D}

% ---- Listings (code blocks) ------------------------------------------------
\lstset{
  basicstyle=\ttfamily\small,
  backgroundcolor=\color{codebg},
  frame=single,
  framesep=6pt,
  rulecolor=\color{codebg},
  breaklines=true,
  columns=fullflexible,
  keepspaces=true,
  xleftmargin=4pt,
  xrightmargin=4pt,
  aboveskip=8pt,
  belowskip=8pt,
}

% ---- Headings --------------------------------------------------------------
\titleformat{\section}{\Large\bfseries\color{eldrunblue}}{\thesection}{0.6em}{}
\titleformat{\subsection}{\large\bfseries\color{eldrungray}}{\thesubsection}{0.6em}{}

\pagestyle{fancy}
\fancyhf{}
\rhead{\small\color{eldrungray}Eldrun}
\lhead{\small\color{eldrungray}Detailed README}
\rfoot{\small\thepage}

% ---- Helper macros ---------------------------------------------------------
\newcommand{\done}{\textcolor{donegreen}{\textbf{Done}}}
\newcommand{\open}{\textcolor{openorange}{\textbf{Open}}}
\newcommand{\code}[1]{\texttt{#1}}

\setlist{itemsep=2pt,topsep=3pt}

\title{\vspace{-1.2cm}\textbf{\color{eldrunblue}Eldrun}\\[0.3em]
\large A project-centric desktop workspace for AI-assisted development}
\author{Detailed technical README}
\date{Version 0.1.0 \quad\textbullet\quad Tauri~2 + React + TypeScript}

\begin{document}

\maketitle
\thispagestyle{fancy}

\begin{abstract}
\noindent
\textbf{The problem:} when you develop several projects at once, every project's
windows --- browsers, terminals, file managers, docs, agents --- pile up on one
desktop. Switching from project~A to project~B means hunting through dozens of
windows for the handful that belong to where you are going, and losing the rest
in the noise.

\smallskip
\noindent
\textbf{Eldrun's model:} you do not open applications, you open \emph{projects}.
Selecting a project swaps the whole desktop to that project's context --- its
windows come forward, the previous project's windows are parked out of the way,
the downloads folder and default-app mappings re-route, and time tracking
switches. One project visible at a time, everything else cleanly out of sight.
\end{abstract}

\tableofcontents
\vspace{0.6cm}

% ===========================================================================
\section{Concept and Vision}
% ===========================================================================

Inside each project, Eldrun is an operational cockpit: a root control terminal
for managing the workspace, one or more agent terminals scoped to the project
(Claude, Codex, Gemini, or a local Ollama model), a persistent project bar, a
hover-revealed file panel, and cross-project app controls that stay available as
you move between projects.

Eldrun is a \emph{project-centric desktop layer}, not just an app that launches
or embeds other apps. The user-facing model is:

\begin{quote}
\itshape Select a project $\rightarrow$ Eldrun restores that project's working
context.
\end{quote}

The core product is the window/workspace layer: projects own their windows and
desktop context, and switching projects swaps that context as one unit. The
agent terminals, file panel, and app launcher ride on top of that layer. That
context should eventually include terminals, files, apps, windows, Git state,
notes, AI/task metadata, layout, and workflow state.

The current implementation targets Linux (X11 and KDE Wayland) because that
provides the fastest path to reliable window control. The longer-term direction
is a stable Eldrun core with desktop/compositor backends for Cinnamon~X11,
KDE/KWin, Hyprland, GNOME Shell, i3, Sway, and other Wayland environments.

\subsection{How Eldrun compares}

\begin{description}[leftmargin=1.2em,style=nextline]
  \item[Agent orchestrators] (Vibe Kanban, Conductor, Claude Squad, the Claude
  Code desktop app) manage agent \emph{processes inside a repo} --- task
  delegation, git worktrees, diff review, merge flow. They have no notion of
  your desktop: they will not move your windows, re-route downloads, or switch
  default apps when you change focus. Eldrun is \emph{complementary} --- you can
  run one inside an Eldrun project terminal.
  \item[Manual approaches] cover only one slice each: KDE Activities and one
  virtual desktop per project handle windows but have no project model and no
  restore; tmux and scripts like \code{workon} restore terminal layouts but
  ignore everything outside the terminal.
\end{description}

Eldrun occupies the gap none of them fill: project ownership of \emph{windows
and desktop context}, with agent terminals built in.

\subsection{Layout sketch}

\begin{lstlisting}
+------------------------------------------------------------------+
| status/network    project switcher pills    app controls         |
+------------------------------------------------------------------+
| global cross-project app toolbar (hover to reveal)               |
+------------------------------------------------------------------+
|                                                                  |
| tiling subwindows: each owns its own tab bar (terminal / agent / | right panel
| file viewer), split L/R/T/B by dragging tabs onto pane edges     | (hover or
|                                                                  |  pin open)
+------------------------------------------------------------------+
\end{lstlisting}

% ===========================================================================
\section{Technology Stack}
% ===========================================================================

\begin{longtable}{@{}p{4.2cm}p{10.5cm}@{}}
\toprule
\textbf{Layer} & \textbf{Components} \\
\midrule
\endhead
Frontend        & React~18, TypeScript, Vite, Tailwind CSS, Zustand \\
Terminal UI     & xterm.js (\code{@xterm/xterm}, \code{@xterm/addon-fit}, \code{@xterm/addon-web-links}) \\
Backend         & Rust, Tauri~v2 \\
PTY             & \code{portable-pty} crate \\
Workspace       & \code{zbus} (DBus) and \code{xcb} (X11) --- Linux only \\
\bottomrule
\end{longtable}

\subsection{Requirements}

\begin{itemize}
  \item Linux desktop (X11 or KDE Wayland)
  \item Rust toolchain (\code{rustup}) and Node 18+
  \item \code{sshfs} + FUSE (optional, only for remote/SSH projects)
\end{itemize}

\begin{lstlisting}[language=bash]
# Install Rust
curl --proto '=https' --tlsv1.2 -sSf https://sh.rustup.rs | sh

# Tauri system dependencies (Debian / Ubuntu)
sudo apt install libwebkit2gtk-4.1-dev libssl-dev libgtk-3-dev \
    libayatana-appindicator3-dev librsvg2-dev

# Install JS deps
npm install
\end{lstlisting}

\subsection{Run}

\begin{lstlisting}[language=bash]
./start-eldrun-tauri.sh          # packaged app
npm run tauri dev                # dev build with hot-reload
\end{lstlisting}

The desktop launchers are \code{Eldrun.desktop} (packaged) and
\code{EldrunHotReload.desktop} (hot reload). They already point at this
checkout's scripts:

\begin{lstlisting}[language=bash]
cp Eldrun*.desktop ~/.local/share/applications/
update-desktop-database ~/.local/share/applications/
\end{lstlisting}

% ===========================================================================
\section{Main Features}
% ===========================================================================

\subsection{Project desktop (the differentiator)}

\begin{description}[leftmargin=1.2em,style=nextline]
  \item[Workspace management] X11 two-desktop parking model and KDE Wayland
  per-project virtual desktop model; global app windows stay visible across all
  project switches.
  \item[External window tracking] file opens use \code{xdg-open}; launched
  windows are tracked by PID and shown in the right panel instead of embedded in
  the UI.
  \item[Downloads routing] \code{\textasciitilde/eldrun/downloads} symlink always
  points to the active project's \code{tmp/downloads/}; Firefox and Chromium
  preferences are updated automatically.
  \item[Default app mapping] file extensions use per-project overrides, global
  defaults, system MIME defaults, or a manual ``Open With'' picker.
  \item[Time tracking] Eldrun records active project sessions and shows today's
  elapsed time on project pills.
\end{description}

\subsection{Project cockpit}

\begin{description}[leftmargin=1.2em,style=nextline]
  \item[Agent-terminal orchestration] create Claude, Codex, Gemini, or plain
  shell tabs from the tab bar; create local Ollama-backed Vibe tabs from
  installed models; rename, close, and reorder them by drag and drop. Tab layout
  is persisted per project.
  \item[Tiling subwindows] the center panel is a tiling layout --- drag a tab
  onto another subwindow's left/right/top/bottom edge to split that direction
  into a new pane, or onto its center to move the tab in. Splits resize with
  draggable dividers, each subwindow keeps its own tab bar, and the whole tree
  is persisted per project.
  \item[Root control terminal] opens in \code{\textasciitilde/eldrun/root/} with
  workspace-level context files.
  \item[Project terminals] each active project gets a PTY tab scoped to its
  directory, with best-effort project-local XDG sandbox paths.
  \item[Project creation and import] the \code{+} button creates a new
  git-backed project or imports an existing directory (keep in place, copy, or
  move).
  \item[Project switcher] search, switch, and close projects; a running-task
  indicator spins on pills with live terminal output (even backgrounded
  projects); hover over a pill to see the project path, status, today's active
  time, and live CPU\%.
  \item[Right file panel] browse, open, create, rename, delete, and reveal
  project files, with a breadcrumb trail and per-file git status markers
  (modified, untracked, staged, committed-but-unpushed, ignored). A ``Git'' view
  shows the current branch, clickable branch pills for checkout, and a commit
  list whose entries open an editable commit-message window (amend HEAD,
  agent-generated messages, or checkout). The panel can be pinned open.
  \item[Global app toolbar] cross-project roles (Browser, Mail, Calendar, File
  Manager, Password Manager, Notes, Screenshot, \dots) with launch-or-raise and
  icon resolution.
  \item[Ollama model management] the Settings Ollama panel shows installed
  models, running CPU/GPU state, parameter and quantization details, plus
  catalog install, update, unload, and delete controls.
  \item[Hover-revealed panels] the global app bar and right file panel appear on
  pointer hover and disappear when the pointer leaves, keeping the center
  terminal unobstructed; the right panel can also be pinned permanently open.
\end{description}

\subsection{Remote (SSH) projects}

Optionally point a project at a remote host. Enter an SSH address
(\code{user@host[:port]}), connect, and browse the remote filesystem in-app to
pick the project root. Eldrun \code{sshfs}-mounts it locally so the file tree,
terminal cwd, and git work unchanged. Terminal and agent tabs run \textbf{on the
remote host} over \code{ssh -tt} (multiplexed over a ControlMaster socket), with
the agent CLI auto-detected/bootstrapped on the remote and authenticated with
the remote's own login. VPN-gated hosts bring up an OpenVPN tunnel first. Auth
uses your existing SSH setup (keys / agent / \code{\textasciitilde/.ssh/config},
\code{BatchMode}); requires \code{sshfs}/FUSE on the local machine.

\subsection{In-app file viewers}

Drag a file from the tree onto a subwindow's tab bar to open it in a tab. The
built-in viewers, by type:

\begin{longtable}{@{}p{2.8cm}p{11.6cm}@{}}
\toprule
\textbf{Type} & \textbf{Viewer} \\
\midrule
\endhead
Text / code &
  (\code{.txt}, \code{.json}, \code{.py}, \code{.rs}, \code{.ts}, \code{.svg},
  \code{.bib}, many more, plus extensionless files like \code{Dockerfile}): an
  editable code editor with a line-number gutter, syntax highlighting,
  Tab/Shift+Tab indent, undo/redo (\code{Ctrl+Z} / \code{Ctrl+Shift+Z}), and a
  save icon (\code{Ctrl+S}). Auto-reloads on disk change (non-destructive
  Reload / Keep-mine banner if you have unsaved edits); offers opt-in local
  autocomplete. \\
\addlinespace
Markdown &
  (\code{.md}, \code{.markdown}, \code{.mdx}): rendered preview with an
  Edit/Preview toggle; links to local files read as clickable. \\
\addlinespace
LaTeX &
  (\code{.tex}): the code editor plus, when a TeX engine is on \code{PATH}, a
  compile action with compiler options (output folder, extra engine flags ---
  shell-escape is always stripped). A successful compile opens the PDF in its
  own tab (reused/refreshed on recompile); \code{Ctrl}/\code{Cmd}+Click follows
  \code{\textbackslash input\{\dots\}} /
  \code{\textbackslash includegraphics\{\dots\}} references. \\
\addlinespace
Images &
  (\code{.png}, \code{.jpg}, \code{.gif}, \code{.webp}, \dots): zoom (to the
  cursor) / pan viewer; the image is also draggable out as an OS drop source. \\
\addlinespace
PDF &
  (\code{.pdf}): rendered with a themed zoom toolbar. \\
\bottomrule
\end{longtable}

Office / spreadsheet formats (\code{.odt}, \code{.xlsx}, \code{.docx}, \dots)
and any other type open in their external default app. Native-viewer behaviour
is configured per file type under \textbf{Settings $\rightarrow$ Native
Viewers}. The text/LaTeX/Markdown editors carry an \code{A-}/\code{A+}
text-size control (or \code{Ctrl} +/-, \code{Ctrl+0} to reset), persisted per
file type.

\subsection{Local autocomplete (opt-in, private)}

In the editable text/LaTeX/markdown viewers, \code{Ctrl+Space} requests a single
completion from a \textbf{local Ollama} model (\code{Tab} accepts, \code{Esc}
dismisses). It is OFF by default and per file type; nothing is sent anywhere
unless you enable it, and if Ollama is not running it fails silently --- no
remote calls, ever.

\subsection{Platform and packaging}

\begin{description}[leftmargin=1.2em,style=nextline]
  \item[Network indicator] probes connectivity and shows online/offline plus
  wired or wireless state.
  \item[Keyboard shortcuts] Eldrun opens fullscreen by default; \code{F11}
  toggles fullscreen; \code{Super} toggles all panels.
  \item[Crash logging] Rust panic hook appends to
  \code{\textasciitilde/.local/share/eldrun/crash.log}.
  \item[Packaging] Debian \code{.deb} and AppImage targets.
\end{description}

% ===========================================================================
\section{Agent Support}
% ===========================================================================

Eldrun launches agents in xterm.js PTY tabs. The active agent command
(\code{claude}, \code{codex}, \code{gemini}, or \code{vibe}) is set in Settings.
If the configured command is not in \code{\$PATH}, Eldrun falls back to the
system shell.

\begin{longtable}{@{}p{4.6cm}p{1.6cm}p{1.4cm}p{6.2cm}@{}}
\toprule
\textbf{Agent} & \textbf{Integ.} & \textbf{Tested} & \textbf{Notes} \\
\midrule
\endhead
Claude (\code{claude})       & Yes & Yes     & Default agent command. Full tab lifecycle, layout persistence, project-scoped sandbox env. \\
Codex (\code{codex})         & Yes & Yes     & Selectable as default agent. Same tab lifecycle as Claude. \\
Gemini (\code{gemini})       & Yes & Yes     & Selectable as default agent. Same lifecycle. \\
Vibe (\code{vibe})           & Yes & No      & Selectable; same tab lifecycle. \\
Ollama via Vibe              & Yes & Partial & Installed Ollama models appear under Local Agents; each local tab gets an isolated per-model \code{VIBE\_HOME}. \\
Shell                        & Yes & Yes     & Plain interactive shell tab in the project directory. \\
Mistral CLI                  & No  & No      & Not integrated; usable in a plain shell tab. \\
Qwen CLI                     & No  & No      & Not integrated. \\
Grok CLI                     & No  & No      & Not integrated. \\
\bottomrule
\end{longtable}

Project-bound terminals receive a best-effort project sandbox: the child process
runs in the project directory with project-local XDG config, cache, data, state,
and temp locations under \code{<project>/.eldrun/sandbox/}. The root
orchestration terminal keeps the normal workspace environment.

\subsection{Session resume}

Claude and Codex tabs that carry a session id are persisted across restarts and
respawned with their prior conversation. Eldrun installs a \code{SessionStart}
hook (into \code{\textasciitilde/.claude/settings.json} and
\code{\textasciitilde/.codex/config.toml}) that records each tab's live session
id keyed by an \code{ELDRUN\_TAB\_UID} env var, so resume follows the live
session even across a \code{/clear}.

\begin{itemize}
  \item For Claude the key is its launch id (\code{--session-id}); the hook
  records the live \code{session\_id} and \code{resolve\_claude\_session} emits
  \code{claude --resume <live-id>}.
  \item Codex mints its own id, so the key is a separate per-tab uuid and the
  backend injects \code{codex resume <live-id>}.
  \item Codex caveat: user-level Codex hooks need a one-time trust (\code{/hooks}
  in Codex) before they fire.
  \item Gemini and Vibe tabs are still dropped on restart.
\end{itemize}

% ===========================================================================
\section{Platform Support}
% ===========================================================================

\begin{longtable}{@{}p{4.6cm}p{3.4cm}p{6.4cm}@{}}
\toprule
\textbf{Platform} & \textbf{Status} & \textbf{Notes} \\
\midrule
\endhead
Linux --- X11          & Yes                 & Two-desktop workspace parking model (EWMH/xcb). Primary development target. \\
Linux --- KDE Wayland  & Yes                 & Per-project virtual desktop via KWin DBus scripting. KDE~5 and KDE~6. \\
Linux --- other Wayland& Partial             & Null backend (no workspace switching, no sticky windows). Terminal and file management work. \\
Windows                & Experimental shell  & Null workspace backend. No native window/default-app/download integrations. \\
macOS                  & Experimental shell  & Null workspace backend. No native window/default-app/download integrations. \\
\bottomrule
\end{longtable}

% ===========================================================================
\section{Project Storage}
% ===========================================================================

Managed projects live under
\code{\textasciitilde/eldrun/projects/<sanitized-name>/}. Imported projects can
also be registered in place. Remote (SSH) projects are sshfs-mounted under
\code{\textasciitilde/.local/share/eldrun/mounts/<project-id>/}.

\noindent Global Eldrun state lives in
\code{\textasciitilde/.local/share/eldrun/}:

\begin{longtable}{@{}p{4.8cm}p{9.8cm}@{}}
\toprule
\textbf{File} & \textbf{Contents} \\
\midrule
\endhead
\code{projects.json}        & Lightweight index: project id, name, status, ordering, path to each project's local metadata. \\
\code{settings.json}        & Default agent command, theme, workspace-management setting, global app registry, preferences. \\
\code{default\_apps.json}   & Global file-extension to application command map. \\
\code{time\_log.json}, \code{active\_session.json} & Session time tracking. \\
\code{boxes.json}           & Project boxes / meta-project grouping (members, folder, relations). \\
\code{vibe\_local/<alias>/config.toml} & Isolated Vibe configuration per local Ollama model tab. \\
\bottomrule
\end{longtable}

Project-local state lives in each project's \code{project.json}, alongside
scaffolded files such as \code{AGENTS.md}, \code{CLAUDE.md}, \code{GEMINI.md},
\code{.claude/settings.json}, \code{.gitignore}, \code{TODO.md},
\code{ROADMAP.md}, \code{STATUS.md}, and \code{README.md}. The per-project tab
layout (\code{tab\_layout}/\code{tab\_groups}) is stored here: shell/files tabs
always restore on relaunch; resumable Claude/Codex agent tabs (carrying a
\code{sessionId}) are restored; other agent tabs are dropped.

% ===========================================================================
\section{Source Map}
% ===========================================================================

\subsection{Frontend (\code{src/})}

\begin{longtable}{@{}p{6.2cm}p{8.4cm}@{}}
\toprule
\textbf{File} & \textbf{Purpose} \\
\midrule
\endhead
\code{App.tsx}                       & Root component, theme injection, global key handlers. \\
\code{main.tsx}                      & React entry point. \\
\code{layout/AppShell.tsx}           & Top-level layout: header, center, bottom, right panel wiring. \\
\code{layout/HeaderBar.tsx}          & Window drag handle, close/minimize/maximize buttons. \\
\code{layout/GlobalAppBar.tsx}       & Global toolbar (app launcher, shortcuts). \\
\code{layout/CenterPanel.tsx}        & Terminal stack and tab bar. \\
\code{layout/ProjectSwitcher.tsx}    & Project pill switcher. \\
\code{layout/RightPanel.tsx}         & File tree overlay panel. \\
\code{projects/ProjectPill.tsx}      & Individual project tab pill. \\
\code{terminal/TerminalView.tsx}     & xterm.js terminal wrapper. \\
\code{stores/projects.ts}            & Zustand store: project list, active project, CRUD. \\
\code{stores/tabs.ts}                & Zustand store for center-panel tabs. \\
\code{stores/settings.ts}            & Zustand store for app settings. \\
\code{stores/windows.ts}             & Zustand store for embedded app windows. \\
\code{stores/boxes.ts}               & Zustand store for project boxes. \\
\code{types/index.ts}                & Shared TypeScript types. \\
\bottomrule
\end{longtable}

\subsection{Backend (\code{src-tauri/src/})}

\begin{longtable}{@{}p{6.2cm}p{8.4cm}@{}}
\toprule
\textbf{File} & \textbf{Purpose} \\
\midrule
\endhead
\code{main.rs}                       & Tauri app entry point, plugin registration. \\
\code{lib.rs}                        & Command registration, app setup. \\
\code{storage.rs}                    & JSON persistence helpers. \\
\code{schema/}                       & Serde structs mirroring JSON state files. \\
\code{schema/boxes.rs}               & \code{ProjectBox}/\code{BoxRelation} model. \\
\code{platform/x11.rs}               & X11 workspace / window management via xlib. \\
\code{platform/wayland\_kde.rs}      & KDE Wayland backend via KWin scripting + DBus. \\
\code{platform/null.rs}              & No-op platform fallback. \\
\code{terminal/}                     & PTY management and terminal I/O commands. \\
\code{commands/}                     & Tauri command handlers exposed to the frontend. \\
\code{commands/ssh.rs}               & SSH commands for remote projects. \\
\code{commands/boxes.rs}             & Box CRUD commands. \\
\code{services/ssh\_mount.rs}        & sshfs mount lifecycle for remote projects. \\
\code{services/ssh\_exec.rs}         & Spawn-rewrite to run tabs on the remote over \code{ssh -tt}. \\
\code{services/openvpn.rs}           & OpenVPN tunnel lifecycle for VPN-gated hosts. \\
\code{services/agent\_session.rs}    & SessionStart hook install + live-id store. \\
\code{services/project\_runtime.rs}  & Project switching coordination. \\
\bottomrule
\end{longtable}

% ===========================================================================
\section{Development Workflow}
% ===========================================================================

\begin{enumerate}
  \item Edit files under \code{src/} (frontend) or \code{src-tauri/src/}
  (backend).
  \item Type-check frontend: \code{npx tsc --noEmit}
  \item Run Rust tests: \code{cargo test --manifest-path src-tauri/Cargo.toml}
  \item Every push should produce a fresh packaged artifact from
  \code{.github/workflows/ci-cd.yml}; use \code{npm run package} locally to
  install the same release build. GitHub Releases are only published for
  \code{v0.<minor>.0} tags.
\end{enumerate}

\noindent\textbf{Hot-reload note.} Frontend (\code{src/}) edits hot-reload in a
running instance --- no restart needed. Only backend (\code{src-tauri/}) changes
require a rebuild/restart. Do not open a second Eldrun instance: it can corrupt
workspace state.

% ===========================================================================
\section{Roadmap and Open Work}
% ===========================================================================

Work items are tracked in \code{docs/open\_ideas.md} with globally stable
numbers (1--65) collected into lettered groups. Status is tracked on three axes:
\done{} (code-complete / compiles), \textbf{Automated} (a vitest or
\code{cargo test} exercises it), and \textbf{Manual} (runtime QA in a live
Eldrun). A feature is fully tested only when both test boxes are ticked ---
currently no manual box is ticked yet.

\begin{longtable}{@{}p{1.6cm}p{7.2cm}p{5.4cm}@{}}
\toprule
\textbf{Group} & \textbf{Theme} & \textbf{State} \\
\midrule
\endhead
A  & Project boxes / meta-project grouping (\#13, \#41) & \done{} (Phase 1 + Phase 2 groundwork); merged multi-root tree \& relation surfacing deferred \\
C  & Workspace switching / platform stability (\#15--\#19) & \open{} --- X11 parking races, z-order, KDE i3-mode, cross-platform QA \\
E  & Git worktree support (\#23) & \open{} --- net-new \\
F  & Session restore (\#24, \#39) & Claude \& Codex per-tab resume \done{}; Gemini/Vibe \open{} \\
G  & Remote / SSH \& containerized projects (\#28, \#38) & SSH \done{} (runtime QA + 28c hardening pending); Docker \open{} \\
H  & Cross-platform Windows \& macOS (\#30, \#31) & \open{} --- builds compile; runtime QA + native window tracking pending \\
I  & Backend runtime follow-ups (\#32) & \open{} --- PTY resurrection, transcript storage, durable window metadata \\
J  & Web \& mail surfaces (in-app mail/browser, routing) & \open{} \\
\bottomrule
\end{longtable}

\subsection{Notable open items}

\begin{description}[leftmargin=1.2em,style=nextline]
  \item[Project box containers (\#41)] open a box as a single \emph{merged}
  workspace spanning its member projects: a multi-root right-panel file tree, a
  \code{\textasciitilde/eldrun/boxes/<name>/} box folder as the cwd for box
  agents, agent tabs seeded with each member's \code{CLAUDE.md}/\code{AGENTS.md}
  hints, and directed inter-project relation edges (``a change in A may
  influence B'') tied into git-status markers. Schema groundwork is in place;
  Phases 3--4 deferred.
  \item[Docker/Podman projects (\#38)] run a project's terminal/agent tabs via
  \code{docker exec} into a container with the project directory bind-mounted;
  mirrors the SSH two-service split (\code{docker\_runtime.rs} +
  \code{docker\_exec.rs}). Phase~1 local, Phase~2 remote-over-SSH.
  \item[SSH hardening (\#28c)] outstanding correctness/security items: remote
  command injection in the browse commands, remote-agent resume ordering bug,
  serializing per-id mounts, env-key validation, VPN activation ordering,
  password-auth half-state, and stale-handle recovery.
  \item[X11 robustness (\#15--\#17)] reliably park opened windows on the hidden
  desktop on switch, harden the two-desktop model against races, and preserve
  per-window z-order across switches.
\end{description}

% ===========================================================================
\section{Current Limits}
% ===========================================================================

\begin{itemize}
  \item Live X11 window embedding (frameless reparenting of an external app into
  a tab) is not yet implemented; files render in built-in viewers where
  available, otherwise open via \code{xdg-open} and are tracked as external
  windows.
  \item KDE Wayland workspace management needs live-session QA.
  \item Tab layout is persisted per project; shell, file-viewer, and resumable
  Claude/Codex agent tabs are restored on relaunch, but other agent tabs
  (Gemini, Vibe) and live PTY scrollback are not.
  \item Non-KDE Wayland compositors fall back to the null backend.
\end{itemize}

\vfill
\begin{center}
\small\color{eldrungray}
See \code{docs/VISION.md} for the full strategy and platform rationale, and
\code{DOCUMENTATION.md} for detailed architecture, data schemas, and behavior
notes.
\end{center}

\end{document}
